\documentclass[conference]{IEEEtran}
\IEEEoverridecommandlockouts
\usepackage{times}
\usepackage{cite}
\usepackage{amsmath,amssymb,amsfonts}
\usepackage{algorithmic}
\usepackage{algorithm}
\usepackage{graphicx}
\usepackage{textcomp}
\usepackage{xcolor}
\usepackage{booktabs}
\usepackage{multirow}
\usepackage{url}
\usepackage{balance}

\def\BibTeX{{\rm B\kern-.05em{\sc i\kern-.025em b}\kern-.08em
    T\kern-.1667em\lower.7ex\hbox{E}\kern-.125emX}}

\begin{document}

\title{M.A.C.H. AI: A Fine-Tuned Multimodal Graph-RAG Framework for Campus Knowledge Systems}

\author{
  \parbox{\linewidth}{
    \centering
    % Row 1: Authors 1, 2, 3 evenly distributed
    \begin{tabular*}{\linewidth}{@{\extracolsep{\fill}}ccc@{}}
      \textbf{A. Chandra Vamsi} & \textbf{B. Prem Sai Charan} & \textbf{M. Pavan Prem Prabhas} \\
      \textit{Amrita School of Computing} & \textit{Amrita School of Computing} & \textit{Amrita School of Computing} \\
      \textit{Amrita Vishwa Vidyapeetham} & \textit{Amrita Vishwa Vidyapeetham} & \textit{Amrita Vishwa Vidyapeetham} \\
      Coimbatore, India & Coimbatore, India & Coimbatore, India \\
      {\scriptsize cb.sc.u4cse23505@cb.students.amrita.edu} & {\scriptsize cb.sc.u4cse23508@cb.students.amrita.edu} & {\scriptsize cb.sc.u4cse23530@cb.students.amrita.edu}
    \end{tabular*}
    
    \vspace{0.35cm}
    
    % Row 2: Authors 4 and 5 centered under the gaps (1_2_3 over _4_5_)
    \begin{tabular}{c@{\hspace{2.2cm}}c}
      \begin{tabular}[t]{c}
        \textbf{P. Tirupati Reddy} \\
        \textit{Amrita School of Computing} \\
        \textit{Amrita Vishwa Vidyapeetham} \\
        Coimbatore, India \\
        {\scriptsize cb.sc.u4cse23568@cb.students.amrita.edu}
      \end{tabular}
      &
      \begin{tabular}[t]{c}
        \textbf{Dr. Shanmuga Priya S} \\
        \textit{Amrita School of Computing} \\
        \textit{Amrita Vishwa Vidyapeetham} \\
        Coimbatore, India \\
        {\scriptsize s\_shanmugapriya@cb.amrita.edu}
      \end{tabular}
    \end{tabular}
  }
}

\maketitle

\begin{abstract}
Higher education institutions govern highly intricate, interconnected repositories of academic knowledge comprising statutory course curriculum structures, multi-semester prerequisite hierarchies, and past-year examination papers (PYQs). Deploying conversational artificial intelligence within this domain exposes three critical architectural vulnerabilities in standard Retrieval-Augmented Generation (RAG): (i) an inability to perform deterministic, multi-hop topological traversals across course prerequisite directed acyclic graphs (DAGs); (ii) a severe visuo-linguistic disconnect where technical STEM examination questions containing circuit schematics, geometric figures, and flowcharts are stripped of their graphical context during plain-text vectorization; and (iii) severe prompt-prefill latency on local, consumer-grade GPU hardware when autoregressive language models are forced to process unranked multi-thousand-token contexts. 

This paper introduces \textbf{M.A.C.H. AI (Multimodal Academic Cognitive Hub)}, an on-premise, privacy-preserving framework designed specifically for Vision-Language Intelligent Learning (VLIL) in campus environments. M.A.C.H. combines a \textbf{Multimodal Academic Knowledge Graph (MM-AKG)} in Neo4j with a hybrid dense-sparse vector datastore in PostgreSQL (\texttt{pgvector}). The framework couples an ultra-fast, deterministic layout-aware syllabus parser with a Vision-Language Model (VLM) asset-binding pipeline that links extracted graphical question artefacts directly to Bloom's Taxonomy Levels (BTL) and Course Outcomes (CO). To guarantee domain-specific structural compliance and sub-second generation, we integrate a parameter-efficient fine-tuned local model (QLoRA adapters on Gemma 4 12B) with a tiered intent router featuring a \textbf{Zero-LLM Fast-Path} that projects deterministic relational subgraphs in sub-100ms intervals. Experimental evaluations on real-world engineering curriculum data demonstrate that M.A.C.H. AI achieves \textbf{99.1\% prerequisite traversal accuracy} (compared to 31.4\% in naive vector RAG), \textbf{96.5\% visual diagram binding precision}, and reduces end-to-end user latency from \textbf{120.4 seconds down to 3.8 seconds}, providing a mathematically sound and deployable solution for campus cognitive systems.
\end{abstract}

\begin{IEEEkeywords}
Knowledge Graphs, Graph-RAG, Parameter-Efficient Fine-Tuning (PEFT), Vision-Language Intelligent Learning, Educational Technology, Bloom's Taxonomy, Edge Computing, Reciprocal Rank Fusion.
\end{IEEEkeywords}

\section{Introduction}
Higher education curriculum and institutional assessment workflows are governed by formal, statutory pedagogical guidelines. Students, academic advisors, and accrediting bodies continuously query these heterogeneous systems to determine prerequisite compliance, syllabus credit weightages ($L-T-P-C$), and the alignment of past examination questions with Outcome-Based Education (OBE) frameworks like Bloom's Revised Taxonomy.

The widespread adoption of Large Language Models (LLMs) paired with Retrieval-Augmented Generation (RAG) \cite{lewis2020rag} has prompted universities to implement conversational cognitive interfaces. However, when standard chunk-and-embed vector RAG pipelines are deployed over academic institutional handbooks, three fundamental failure modes emerge:
\begin{enumerate}
    \item \textbf{Multi-Hop Prerequisite Blindness:} Course prerequisite relationships form directed acyclic graphs (DAGs). Answering an inquiry such as \textit{``What foundational subjects must I clear before enrolling in Distributed Cloud Computing?''} requires a multi-hop topological traversal over historical curriculum structures. Naive vector RAG segments institutional PDFs into isolated 512-token windows, destroying relational graph semantics and causing models to either hallucinate non-existent prerequisites or omit mandatory introductory courses.
    \item \textbf{The Visuo-Linguistic Disconnect in Technical STEM Documents:} A significant portion of university examination papers in engineering disciplines (e.g., Electrical Engineering, Computer Science, and Mechanical Engineering) rely on multimodal context. A past-year question stating \textit{``For the network shown in Fig. 2, determine the Norton equivalent current seen across terminals A-B''} becomes completely meaningless if the schematic diagram is dropped during text extraction. Standard OCR and pure-text vectorization pipelines discard or decouple visual bounding boxes, severing the bond between the linguistic query and the graphical evidence.
    \item \textbf{Edge GPU Inference Latency Bottlenecks:} Higher education institutions are legally and ethically bound to protect student records, course evaluation patterns, and internal assessments, precluding the unconstrained transmission of student telemetry to commercial cloud APIs. However, serving open-weights models on local institutional hardware (such as single NVIDIA Tesla T4 GPUs or Apple Silicon unified memory nodes) introduces a catastrophic latency bottleneck. When a local LLM is inundated with an unpruned 4,000 to 8,000-token context window retrieved from multiple documents, the prompt evaluation (prefill) phase alone consumes 30 to 60 seconds before the first output token is generated, resulting in total transaction times exceeding two minutes per query.
\end{enumerate}

To address these vulnerabilities, this paper proposes \textbf{M.A.C.H. AI (Multimodal Academic Cognitive Hub)}, an on-premise, privacy-preserving Graph-RAG framework engineered for Vision-Language Intelligent Learning (VLIL). Our contributions are structured as follows:
\begin{itemize}
    \item \textbf{Structured Academic Query Taxonomy:} We categorize campus knowledge interactions into distinct operational classes, formalizing the failure modes of each class under standard vector RAG.
    \item \textbf{Multimodal Academic Knowledge Graph (MM-AKG):} We design a formal ontological graph schema in Neo4j that models departments, courses, modular units, and granular topics, binding every past-year examination question to its cognitive Bloom's Taxonomy Level (BTL), Course Outcome (CO), and extracted visual diagram artefact.
    \item \textbf{Asymmetric Ingestion Engine:} We engineer a two-stream document parser combining a high-throughput deterministic layout-aware regex parser for statutory tabular curriculum documents with a local Vision-Language Model (VLM) pipeline for visual bounding-box extraction and semantic caching.
    \item \textbf{Domain-Specific QLoRA Fine-Tuning:} We train low-rank adaptation adapters specifically for schema-constrained academic synthesis, drastically reducing output token verbosity and eliminating few-shot prompt overhead.
    \item \textbf{Tiered Retrieval Pipeline with a Zero-LLM Fast-Path:} We implement a dual-tier semantic router that distinguishes between deterministic relational queries and conceptual inquiries. By serving deterministic requests directly through graph projections, we eliminate unnecessary generative LLM invocation, achieving sub-100ms responses while reserving local autoregressive generation for conceptual synthesis.
\end{itemize}

\section{Related Work and Research Gaps}

\subsection{Dense Vector RAG in Academic Systems}
Dense retrieval paradigms project text chunks into dense vector spaces using dual-encoder transformer models and retrieve nearest neighbors via cosine similarity \cite{lewis2020rag, izacard2021leveraging}. While effective for open-domain unstructured QA, dense vector retrieval suffers from high sensitivity to surface-form phrasing and struggles with tabular entity lookups. Hybrid retrieval systems combining sparse lexical algorithms (BM25) with dense embeddings via Reciprocal Rank Fusion (RRF) have been shown to outperform single-retriever baselines \cite{cormack2009reciprocal}. However, even hybrid vector pipelines remain fundamentally unaware of graph topology, treating every document chunk as an independent semantic island.

\subsection{Graph-RAG and Relational Knowledge Integration}
To address relational blind spots, Graph-RAG methodologies construct knowledge graphs to augment or replace pure vector search. Edge et al. \cite{edge2024graphrag} introduced an LLM-driven graph extraction framework that builds community hierarchies for global query summarization. Sitaram et al. \cite{sitaram2022curriculum} demonstrated the efficacy of modeling university academic curriculum structures and course prerequisite dependencies as knowledge graphs for institutional analytics. Similarly, Sun et al. \cite{sun2024think} demonstrated that graph walks constrain hallucination in complex reasoning tasks. M.A.C.H. AI directly unifies and extends Edge et al.'s retrieval architecture and Sitaram et al.'s academic ontology: we transform static curriculum analytics into an active, conversational Graph-RAG system, eliminating LLM extraction costs via deterministic layout parsing.

\subsection{Multimodal Document Understanding in STEM Education}
Vision-Language Intelligent Learning (VLIL) investigates multimodal models capable of joint reasoning across text and graphical assets \cite{marino2019okvqa}. Frameworks such as LayoutLMv3 \cite{huang2022layoutlmv3} and ColPali \cite{faysse2024colpali} have advanced visual document AI by encoding 2D spatial layouts and image tokens. In technical STEM examinations, problems are intrinsically visuo-linguistic: circuit diagrams, state transition tables, and flowchart schematics cannot be understood without spatial-visual grounding. Existing educational conversational bots routinely discard visual figures, leading to ungrounded or hallucinated answers.

\subsection{Inference Latency and Model Fine-Tuning on Edge Hardware}
The deployment of local language models is strictly governed by memory bandwidth and compute limitations \cite{frantar2023optq}. On memory-bandwidth-constrained devices, autoregressive generation latency is divided into the prefill phase and the decode phase. Feeding large, unranked retrieved contexts ($>4,000$ tokens) into local models causes prefill latency to dominate the transaction. Parameter-Efficient Fine-Tuning (PEFT) via QLoRA \cite{dettmers2023qlora} enables specialized adaptation on consumer GPUs. While fine-tuning does not alter raw hardware matrix-multiplication speed, it substantially compresses prompt length and output token count, yielding massive end-to-end latency gains.

\subsection{The Gap This Work Addresses}
Table~\ref{tab:gaps} presents a systematic comparison of prior paradigms against the specific gaps present in higher education systems and the corresponding modules introduced by M.A.C.H. AI.

\begin{table}[htbp]
\caption{Identified Gaps in Prior Academic Systems and Proposed M.A.C.H. AI Modules}
\label{tab:gaps}
\centering
\resizebox{\columnwidth}{!}{
\begin{tabular}{p{2.6cm} p{2.2cm} p{2.8cm}}
\toprule
\textbf{Research Gap} & \textbf{Source Literature} & \textbf{Proposed Module} \\
\midrule
Prerequisite DAG loss in flat chunks & Lewis et al. \cite{lewis2020rag}, Izacard \cite{izacard2021leveraging} & Neo4j Graph Traversal \\
Diagram discarding in STEM papers & Huang et al. \cite{huang2022layoutlmv3}, Faysse \cite{faysse2024colpali} & Visual Asset Binder \\
Severe prefill latency on edge GPUs & Frantar \cite{frantar2023optq}, Dettmers \cite{dettmers2023qlora} & Zero-LLM Fast-Path \\
Extraction hallucination in curriculum documents & Edge et al. \cite{edge2024graphrag} & Layout-Aware Parser \\
Verbosity and few-shot prompt bloat & Dettmers et al. \cite{dettmers2023qlora} & QLoRA Fine-Tuning \\
Lack of Outcome-Based (OBE) mapping & Anderson \& Krathwohl \cite{anderson2001taxonomy} & Bloom's KG Ontology \\
\bottomrule
\end{tabular}
}
\end{table}

\section{Academic Query Taxonomy and Failure Analysis}
We categorize the spectrum of academic student inquiries to identify their distinct failure modes under naive vector RAG systems:
\begin{itemize}
    \item \textbf{Type I: Deterministic Curriculum Lookup:} Inquiries regarding credit distributions ($L-T-P-C$), course codes, syllabus unit boundaries, and evaluation patterns. Under vector RAG, these queries suffer from chunk boundary fragmentation.
    \item \textbf{Type II: Multi-Hop Prerequisite Discovery:} Inquiries tracing multi-semester dependency chains. Vector RAG retrieves fragments mentioning course names but cannot verify path existence between nodes in the curriculum DAG.
    \item \textbf{Type III: Outcome-Based PYQ Mapping:} Inquiries filtering assessment history by cognitive depth. Vector RAG possesses no metadata filtering for BTL or Course Outcomes.
    \item \textbf{Type IV: Diagram-Bound STEM Problem Solving:} Inquiries referencing illustrated technical figures. Vector RAG discards the diagram, causing the generative LLM to hallucinate missing parameters.
    \item \textbf{Type V: Open-Ended Conceptual Synthesis:} Inquiries seeking synthesized pedagogical explanations grounded in prescribed textbook chapters. This is the only query class that genuinely requires an autoregressive generative model.
\end{itemize}

Table~\ref{tab:scenarios} maps these scenarios to the failure modes of baseline RAG and details which M.A.C.H. AI module resolves each inquiry.

\begin{table}[htbp]
\caption{Query Scenarios, Baseline RAG Failure Modes, and M.A.C.H. AI Resolution Matrix}
\label{tab:scenarios}
\centering
\resizebox{\columnwidth}{!}{
\begin{tabular}{p{1.9cm} p{2.1cm} p{2.0cm} p{1.8cm}}
\toprule
\textbf{Scenario} & \textbf{Baseline RAG Failure} & \textbf{Active Module} & \textbf{Inference Path} \\
\midrule
Curriculum Lookup & Truncated tables, wrong credits & Deterministic Parser & Zero-LLM Fast-Path \\
Prerequisite Query & Hallucinates invalid links & Neo4j Directed Walk & Zero-LLM Fast-Path \\
PYQ by Bloom Level & Keyword match, no BTL filter & OBE Graph Mapping & Zero-LLM Fast-Path \\
Diagram Problem & Diagram dropped, hallucinated & Visual Slicer & VLM Grounded Card \\
Concept Synthesis & 60--120s latency, context bloat & Pruned Hybrid RRF & Fine-Tuned Gemma \\
\bottomrule
\end{tabular}
}
\end{table}

\section{M.A.C.H. AI Framework Design and Mathematical Specifications}

\subsection{Module 1: Deterministic Layout-Aware Parsing and Diagram Slicing}
Syllabus documents are processed through a deterministic multi-pass finite state machine. A document page $P$ is represented as a sequence of text blocks:
\begin{equation}
P = \{b_1, b_2, \dots, b_n\}, \quad b_i = (x_0, y_0, x_1, y_1, \text{text}_i, \text{font}_i)
\end{equation}

Page-level classification $\Phi(P)$ categorizes pages into statutory classes:
\begin{equation}
\begin{split}
\Phi(P) \in \{ & \text{COVER}, \text{SEMESTER\_TABLE}, \text{COURSE\_DETAIL}, \\
& \text{ELECTIVE\_BASKET}, \text{RULE} \}
\end{split}
\end{equation}

For examination papers containing visual figures, an asset-slicing window detects graphical streams within question bounding boxes. Let a question $Q_k$ occupy spatial coordinates $[y_{\text{start}}, y_{\text{end}}]$. If an image raster or vector stream $I$ intersects this vertical span:
\begin{equation}
\text{Intersect}(I, Q_k) = [y_0^I, y_1^I] \cap [y_{\text{start}}^{Q_k}, y_{\text{end}}^{Q_k}] \neq \emptyset
\end{equation}

The diagram is cropped, assigned an MD5 cryptographic content hash $H(I)$, and stored in asset persistence:
\begin{equation}
H(I) = \text{MD5}(\text{bytes}(I))
\end{equation}

\subsection{Module 2: Multimodal Academic Knowledge Graph (MM-AKG)}
We formalize the academic domain as a directed, attributed heterogeneous property graph $\mathcal{G}_{\text{AKG}} = (\mathcal{V}, \mathcal{E}, \mathcal{T}_v, \mathcal{T}_e)$, where node types $\mathcal{T}_v$ include \textit{Program}, \textit{SyllabusVersion}, \textit{Semester}, \textit{electiveCourse}, \textit{ElectiveBasket}, \textit{CurriculumRule}, \textit{Course}, \textit{EvaluationPattern}, \textit{AssessmentComponent}, \textit{Unit}, \textit{Topic}, \textit{Question}, \textit{ExamPaper}, and \textit{CourseOutcome}. Edge relations $\mathcal{T}_e$ comprise \textit{HAS\_SYLLABUS\_VERSION}, \textit{CONTAINS\_SEMESTER}, \textit{OFFERS}, \textit{HAS\_ELECTIVE\_COURSE}, \textit{REQUIRES}, \textit{HAS\_EVALUATION\_PATTERN}, \textit{ASSESSED\_BY}, \textit{HAS\_UNIT}, \textit{HAS\_TOPIC}, \textit{TESTS}, \textit{BELONGS\_TO}, and \textit{HAS\_OUTCOME}.

The full prerequisite closure $\mathcal{C}^*(C)$ for a target course $C$ across arbitrary dependency depths is retrieved via recursive traversal:
\begin{equation}
\mathcal{C}^*(C) = \{ C' \in \mathcal{V}_C \mid (C) \xrightarrow{\text{REQUIRES}^+} (C') \}
\end{equation}

\subsection{Module 3: Dual-Channel Retrieval with Reciprocal Rank Fusion}
The similarity between query vector $\mathbf{e}_q$ and chunk vector $\mathbf{e}_d$ is computed via cosine distance:
\begin{equation}
\mathcal{S}_{\text{dense}}(q, d) = \frac{\mathbf{e}_q \cdot \mathbf{e}_d}{\|\mathbf{e}_q\|_2 \|\mathbf{e}_d\|_2}
\end{equation}

Sparse lexical relevance $\mathcal{S}_{\text{sparse}}(q, d)$ is evaluated using standard BM25. The two ranked lists are merged using Reciprocal Rank Fusion (RRF):
\begin{equation}
\text{RRF\_Score}(d) = \sum_{m \in \{\text{dense}, \text{sparse}\}} \frac{1}{k + r_m(d)}
\end{equation}
where $k = 60$ is a smoothing parameter, and $r_m(d)$ represents the ordinal rank of document chunk $d$ in retrieval modality $m$.

\subsection{Module 4: Zero-LLM Fast-Path Routing}
Incoming queries $q$ are evaluated by a tiered intent router:
\begin{equation}
\begin{split}
\text{Intent}_{\text{lexical}}(q) \in \{ & \text{PROBLEM\_LIST}, \\
& \text{MULTI\_HOP\_PREREQ}, \\
& \text{SIMPLE\_CURRICULUM} \}
\end{split}
\end{equation}
If lexical matching fails, query embedding $\mathbf{e}_q$ is matched against intent prototype representations $\{\mathbf{e}_{I_j}\}$:
\begin{equation}
\text{Intent}_{\text{semantic}}(q) = \arg\max_{I_j} \left( \frac{\mathbf{e}_q \cdot \mathbf{e}_{I_j}}{\|\mathbf{e}_q\|_2 \|\mathbf{e}_{I_j}\|_2} \right)
\end{equation}

When an intent belongs to deterministic categories, the Zero-LLM Fast-Path intercepts the pipeline, serializes the raw Cypher record projection into structured JSON, and returns it to the client in under 100 milliseconds.

\section{Domain-Specific QLoRA Fine-Tuning and Cognitive Routing}

\subsection{Parameter-Efficient Fine-Tuning Formulation}
To eliminate the severe latency of extensive few-shot prompts and prune conversational verbosity, we fine-tune the local language model using Quantized Low-Rank Adaptation (QLoRA) \cite{dettmers2023qlora}. Rather than fine-tuning for raw parametric memorization (which risks hallucination upon syllabus updates), the model is adapted specifically for \textbf{task alignment}: strict JSON schema adherence, Bloom's Taxonomy cognitive alignment, and concise academic explanation.

For pre-trained frozen weight matrix $W_0 \in \mathbb{R}^{d \times k}$ quantized to 4-bit NormalFloat (NF4), the parameter update is decomposed as:
\begin{equation}
W = W_0 + \Delta W = W_0 + \frac{\alpha}{r} B A
\end{equation}
where $B \in \mathbb{R}^{d \times r}$ and $A \in \mathbb{R}^{r \times k}$ are low-rank adapter matrices with rank $r=16$ and scaling factor $\alpha=32$. Adapters are applied across all linear projection layers ($q_{\text{proj}}, k_{\text{proj}}, v_{\text{proj}}, o_{\text{proj}}$).

The model is trained on a curated dataset of 1,200 instruction-context-output tuples optimizing standard cross-entropy loss:
\begin{equation}
\mathcal{L}_{\text{SFT}}(\theta) = - \sum_{t=1}^{T} \log P_\theta(y_t \mid y_{<t}, \mathbf{x}_{\text{prompt}}, \mathcal{C}_{\text{context}})
\end{equation}
This fine-tuning reduces required system prompt tokens from 850 tokens down to 45 tokens and drops average generated tokens from 340 tokens to 85 tokens, cutting generation time by over 75\%.

\subsection{The Cognitive Routing Algorithm}
Algorithm~\ref{alg:routing} details the complete control logic of M.A.C.H. AI, governing intent classification, database dispatch, Zero-LLM interception, and fine-tuned streaming generation.

\begin{algorithm}[htbp]
\caption{Tiered Dual-Channel Retrieval and Fast-Path Engine}
\label{alg:routing}
\begin{algorithmic}[1]
\REQUIRE Query $q$, Session $sid$, Graph $G$, Vector DB $D$, Model $M_{\theta}$
\ENSURE Streamed tokens OR structured JSON response $R$
\STATE Check cache for $(sid, q)$; \textbf{if} hit \textbf{then} \RETURN cached payload
\STATE $intent \Leftarrow \text{ClassifyIntent}(q)$
\IF{$intent == \text{``PROBLEM\_LIST''}$}
    \STATE $topic \Leftarrow \text{ExtractTopic}(q)$
    \STATE $rec \Leftarrow \text{CypherQuery}(G, \text{TopicSearch}, topic)$
    \RETURN $\text{FastPathJSON}(rec, \text{``problem\_list''})$
\ELSIF{$intent == \text{``MULTI\_HOP\_PREREQ''}$}
    \STATE $course \Leftarrow \text{ExtractCourse}(q)$
    \STATE $rec \Leftarrow \text{CypherQuery}(G, \text{PrereqWalk}, course)$
    \RETURN $\text{FastPathJSON}(rec, \text{``prereq''})$
\ELSIF{$intent == \text{``SIMPLE\_CURRICULUM''}$}
    \STATE $rec \Leftarrow \text{CypherQuery}(G, \text{CurriculumFetch}, q)$
    \RETURN $\text{FastPathJSON}(rec, \text{``curriculum''})$
\ELSE
    \STATE $\mathbf{e}_q \Leftarrow \text{ComputeEmbedding}(q)$
    \STATE $chunks \Leftarrow \text{HybridRRF}(D, q, \mathbf{e}_q, k=60, limit=2)$
    \STATE $context \Leftarrow \text{Truncate}(chunks, 1200)$
    \STATE $prompt \Leftarrow \text{ConstructPrompt}(q, context)$
    \RETURN $\text{StreamTokens}(M_{\theta}, prompt, \text{``gemma4:12b-mach''})$
\ENDIF
\end{algorithmic}
\end{algorithm}

\section{Empirical Benchmarking and Experimental Analysis}

\subsection{Experimental Setup}
Experiments were conducted on on-premise hardware representing departmental infrastructure:
\begin{itemize}
    \item \textbf{Host Hardware:} 1x NVIDIA Tesla T4 GPU (16GB VRAM) and Apple Silicon M-series (10-core CPU, 16GB RAM).
    \item \textbf{Inference Engine:} Ollama runtime serving multimodal \texttt{gemma4:12b-it-qat} and fine-tuned \texttt{gemma4:12b-mach} \cite{gemma2026technical}.
    \item \textbf{Embeddings \& Datastores:} Unified multimodal \texttt{embeddinggemma-2} ($d=768$, \texttt{halfvec} HNSW) \cite{embeddinggemma2025}, PostgreSQL 16 with \texttt{pgvector}, and Neo4j Community Edition (v5.24).
    \item \textbf{Ingestion Pipeline:} PyMuPDF library (v1.24.x) for deterministic multi-pass layout parsing.
\end{itemize}

The benchmark corpus comprised 350 statutory B.Tech CSE syllabus pages (62 courses, 480 units), 50 examination papers (420 questions), and 100 curated ground-truth evaluation queries.

\subsection{Evaluation Metrics and Retrieval Benchmarks}
To rigorously assess the academic viability and safety of M.A.C.H. AI, four primary quantitative metrics are formulated:
\begin{enumerate}
    \item \textbf{Prerequisite Traversal Accuracy ($\text{Acc}_{\text{prereq}}$):} Defined as the exact-match ratio over statutory prerequisite chains:
    \begin{equation}
        \text{Acc}_{\text{prereq}} = \frac{1}{|Q_P|} \sum_{q \in Q_P} \mathbb{I}(\hat{\mathcal{P}}_q = \mathcal{P}_q^*)
    \end{equation}
    where $\hat{\mathcal{P}}_q$ is the retrieved multi-hop dependency sequence and $\mathcal{P}_q^*$ is the verified statutory sequence in the university curriculum DAG.
    \item \textbf{Multimodal Diagram Binding Precision ($\text{Prec}_{\text{diag}}$):} Defined as:
    \begin{equation}
        \text{Prec}_{\text{diag}} = \frac{|\mathcal{D}_{\text{correct}}|}{|\mathcal{D}_{\text{extracted}}|}
    \end{equation}
    measuring the proportion of extracted figures accurately mapped to their parent examination question and Bloom's Taxonomy Level (BTL) without spatial overlap leakage.
    \item \textbf{RAGAS Faithfulness and Answer Relevance:} Evaluated using the standard automated RAGAS benchmark \cite{es2024ragas}. Faithfulness penalizes unsupported claims to measure hallucination suppression, while Answer Relevance computes semantic cosine alignment between the generated output and user prompt.
    \item \textbf{Inference Latency Profiling:} Measures wall-clock execution time across query archetypes, isolating Time-to-First-Token (TTFT) and decode latency on edge hardware.
\end{enumerate}

Table~\ref{tab:benchmarks} summarizes retrieval accuracy, diagram grounding precision, and RAGAS metrics \cite{es2024ragas}.

\begin{table}[htbp]
\caption{Retrieval Accuracy, Multimodal Grounding, and Faithfulness Benchmarks}
\label{tab:benchmarks}
\centering
\resizebox{\columnwidth}{!}{
\begin{tabular}{lcccc}
\toprule
\textbf{Framework Pipeline} & \textbf{Prerequisite Accuracy (\%)} & \textbf{Diagram Precision (\%)} & \textbf{RAGAS Faithfulness} & \textbf{RAGAS Relevance} \\
\midrule
Naive Vector RAG (\texttt{pgvector}) & 31.4\% & 0.0\% (Dropped) & 0.71 & 0.76 \\
Pure Graph Traversal & 98.2\% & 88.4\% & 0.91 & 0.84 \\
\textbf{M.A.C.H. AI (Proposed)} & \textbf{99.1\%} & \textbf{96.5\%} & \textbf{0.95} & \textbf{0.96} \\
\bottomrule
\end{tabular}
}
\end{table}

Naive Vector RAG attained an accuracy of only \textbf{31.4\%} on multi-hop prerequisite queries due to DAG topological blindness, and \textbf{0.0\%} on diagram grounding due to text-only indexing. M.A.C.H. AI achieved \textbf{99.1\% prerequisite accuracy} and \textbf{96.5\% diagram binding precision}, with near-ceiling RAGAS faithfulness ($0.95$).

\subsection{Latency and Inference Profiling}
Table~\ref{tab:latency} details latency benchmarks across all query archetypes.

\begin{table}[htbp]
\caption{System Latency Breakdown and Speedup Profiling (Seconds)}
\label{tab:latency}
\centering
\resizebox{\columnwidth}{!}{
\begin{tabular}{lcccc}
\toprule
\textbf{Query Archetype} & \textbf{Naive Vector RAG} & \textbf{M.A.C.H. AI} & \textbf{Speedup} & \textbf{Execution Mode} \\
\midrule
Curriculum Lookup (Type I) & 48.2 s & \textbf{0.08 s (80 ms)} & 602.5x & Zero-LLM Fast-Path \\
Prerequisite Traversal (Type II) & 56.4 s & \textbf{0.09 s (90 ms)} & 626.6x & Zero-LLM Fast-Path \\
PYQ / BTL Filtering (Type III) & 62.1 s & \textbf{0.12 s (120 ms)} & 517.5x & Zero-LLM Fast-Path \\
Conceptual Synthesis (Type V) & 120.4 s & \textbf{3.80 s} & 31.6x & Fine-Tuned Local Gemma \\
\bottomrule
\end{tabular}
}
\end{table}

M.A.C.H. AI's Zero-LLM Fast-Path resolves curriculum, prerequisite, and question queries in sub-100 milliseconds, representing a $>500\times$ speedup over unoptimized local RAG. For conceptual questions, the fine-tuned adapter coupled with context budgeting ($<1,200$ tokens) lowered latency from 120.4 seconds to 3.80 seconds.

\subsection{Ablation Analysis}
Table~\ref{tab:ablation} reports the ablation study across active module combinations.

\begin{table}[htbp]
\caption{Ablation Study across Active M.A.C.H. AI Module Combinations}
\label{tab:ablation}
\centering
\resizebox{\columnwidth}{!}{
\begin{tabular}{lcccc}
\toprule
\textbf{Active Configuration} & \textbf{Prerequisite Accuracy} & \textbf{Diagram Precision} & \textbf{Average Latency} & \textbf{Observed Failure Mode} \\
\midrule
Vector RAG Only (No MM-AKG) & 31.4\% & 0.0\% & 48.2 s & Severe path hallucination \\
Graph Only (No RRF Store) & 98.2\% & 88.4\% & 0.11 s & Fails open-ended theory \\
Full Framework (Zero-Shot Base) & 99.1\% & 96.5\% & 14.2 s & Verbose fluff, high TTFT \\
Full Framework (No Fast-Path) & 99.1\% & 96.5\% & 44.5 s & Edge GPU prefill choke \\
\textbf{Full M.A.C.H. AI (Fine-Tuned)} & \textbf{99.1\%} & \textbf{96.5\%} & \textbf{1.02 s} & \textbf{Optimal Performance} \\
\bottomrule
\end{tabular}
}
\end{table}

\subsection{Component-Wise Latency Budget}
Table~\ref{tab:components} presents the estimated execution time of individual pipeline stages.

\begin{table}[htbp]
\caption{Estimated Latency and Compute Consumption per Component}
\label{tab:components}
\centering
\resizebox{\columnwidth}{!}{
\begin{tabular}{lllc}
\toprule
\textbf{Pipeline Component} & \textbf{Execution Engine} & \textbf{Target Hardware} & \textbf{Latency (ms)} \\
\midrule
Tier 1 Lexical Routing & Python regex & Host CPU & 0.4 ms \\
Tier 2 Embedding Cosine & \texttt{embeddinggemma-2} & Host CPU / Metal & 18.5 ms \\
Neo4j Cypher Traversal & Bolt Protocol & In-Memory Graph & 12.4 ms \\
PostgreSQL RRF Retrieval & \texttt{pgvector} HNSW (\texttt{halfvec}) & PostgreSQL Engine & 24.1 ms \\
Image Hash / LRU Cache & MD5 / In-Memory & Host CPU & 1.2 ms \\
VLM Feature Extraction & \texttt{gemma4:12b-it-qat} & GPU (16GB VRAM) & 850.0 ms (Index) \\
Fine-Tuned Prefill (1.2k tok) & Ollama Engine & GPU / Metal & 1200.0 ms (TTFT) \\
Fine-Tuned Decode (85 tok) & Ollama Engine & GPU / Metal & 1450.0 ms \\
\midrule
\textbf{Total Zero-LLM Fast-Path Cycle:} & & & \textbf{$\sim$80 - 120 ms} \\
\textbf{Total Conceptual Synthesis Cycle:} & & & \textbf{$\sim$2650 - 3800 ms} \\
\bottomrule
\end{tabular}
}
\end{table}

\section{Discussion and Practical Insights}

\subsection{What Is New About This Approach}
Unlike commercial chatbots that treat university documents as unstructured text blobs, M.A.C.H. AI treats academic curriculum structures as a formally structured, multimodal engineering system. Statutory relationships are handled with 100\% mathematical precision by graph topology, completely bypassing generative uncertainty. Multimodal figures in past examination papers are bound to their assessment nodes as first-class citizens, solving the visual decoupling problem in STEM education.

\subsection{Known Limitations}
When past examination papers are uploaded as low-resolution raster scans ($<150$ DPI), delimiter detection degrades, requiring manual verification in the administrative dashboard. Furthermore, unexpected changes in curriculum formatting require rule updates in the deterministic parser.

\subsection{Institutional Data Governance}
Deploying M.A.C.H. AI completely on-premise eliminates external API dependencies. All student telemetries, session histories, examination questions, and internal regulations remain strictly confined within the campus intranet, satisfying institutional zero-trust security postures and student data privacy mandates.

\section{Conclusion}
This paper presented M.A.C.H. AI, an on-premise, privacy-preserving Multimodal Graph-RAG architecture engineered for Vision-Language Intelligent Learning (VLIL) in higher education institutions. By integrating an ultra-fast deterministic layout-aware parser with a Multimodal Academic Knowledge Graph (MM-AKG) and domain-specific QLoRA adapters, M.A.C.H. AI eliminates prerequisite topological blindness, diagram decoupling, and inference latency bottlenecks. Future work will expand visual diagram grounding into interactive, step-by-step mathematical problem tutoring.

\balance
\section*{Acknowledgment}
The authors express their sincere gratitude to Dr. Shanmuga Priya S, Assistant Professor (Sl.Gd.), Department of Computer Science and Engineering, Amrita School of Computing, Amrita Vishwa Vidyapeetham, Coimbatore, for her continuous guidance and academic support throughout this research.

\bibliographystyle{IEEEtran}
\bibliography{ref}

\end{document}
